% s99-bibliography.tex —— 文献表（原文 p.53–54，共 42 条）
% 体例：条目文字保持原文（法文/英文）原样，不译、不删、不改年份与作者；
%       只补 \bibitem 标号与 \textit/\url 排版，与原文 [A1]–[W2] 的字母编号一致。

\begin{thebibliography}{W2}
\srcpage{53}
\bibitem{A1} J. Arthur, \textit{Unipotent automorphic representations: conjectures}, Astérisque 171--172 (1989), p. 13--71, and \textit{Unipotent automorphic representations: global motivation}, in \textit{Automorphic Forms, Shimura varieties, and L-functions}, Academic Press (1990).
\bibitem{A2} J. Arthur, \textit{An introduction to the trace formula}, in [AEK].
\bibitem{A3} J. Arthur, \textit{The Endoscopic Classification of Representations: Orthogonal and Symplectic Groups}, AMS Colloquium Publications, to appear.
\bibitem{AEK} J. Arthur, D. Elwood and R. Kottwitz (eds.), \textit{Harmonic analysis, the trace formula, and Shimura varieties}, American Mathematical Society, p. 1--265 (2005).
\bibitem{BC} A. Borel and B. Casselman (eds.), \textit{Automorphic forms, representations, and L-functions}, Proc. Symp. in Pure Mathematics 33, American Mathematical Society (1977).
\bibitem{BM} A. Borel and G. D. Mostow (eds.), \textit{Algebraic groups and discontinuous subgroups}, Proc. Symp. in Pure Mathematics, American Mathematical Society (1966).
\bibitem{BT} A. Borel et J. Tits, \textit{Groupes réductifs}, Publ. Math. de l'IHES 27 (1965), p. 55--150.
\bibitem{Bo} N. Bourbaki, \textit{Groupes et algèbres de Lie}, Chaps. 4, 5 et 6, Hermann (1968).
\bibitem{D} P. Deligne, \textit{La conjecture de Weil. II}, Publ. Math. de l'IHES 52 (1980), p. 137--252.
\bibitem{DG} R. Y. Donagi et D. Gaitsgory, \textit{The gerbe of Higgs bundles}, Transformation groups 7 (2002), p. 109--154.
\bibitem{FN} E. Frenkel et B. C. Ngô, \textit{Geometrization of trace formulas}, Preprint \url{https://arxiv.org/abs/1004.5323}.
\bibitem{G} Jayce Getz, en préparation.
\bibitem{H} Jochen Heinloth, \textit{Uniformization of G-bundles}, à paraître dans Math. Ann., disponible à \url{http://staff.science.uva.nl/~heinloth/}.
\bibitem{IK} Henryk Iwaniec et Emmanuel Kowalski, \textit{Analytic number theory}, Amer. Math. Soc. Colloquium Publications, vol. 53 (2000).
\bibitem{K1} R. Kottwitz, \textit{Rational conjugacy classes in reductive groups}, Duke Math. J. 49 (1982), p. 785--806.
\bibitem{K2} R. Kottwitz, \textit{Stable trace formula: cuspidal tempered terms}, Duke Math. J. 51 (1984), p. 611--650.
\bibitem{K3} R. Kottwitz, \textit{Tamagawa numbers}, Ann. of Math. 127 (1988), p. 629--646.
\bibitem{LL} J.-P. Labesse et R. Langlands, \textit{L-indistinguishability for SL(2)}, J. Can. de Math. 31 (1979), p. 726--785.
\bibitem{La} S. Lang, \textit{Algebraic number theory}, Addison-Wesley (1968).
\bibitem{L1} R. P. Langlands, \textit{On Artin's L-functions}, Rice University Studies 56 (1970).
\bibitem{L2} R. P. Langlands, \textit{Les débuts d'une formule des traces stable}, Publ. Math. de l'Univ. Paris VII (1981).
\srcpage{54}
\bibitem{L3} R. P. Langlands, \textit{Beyond endoscopy}, \url{http://publications.ias.edu/rpl/}.
\bibitem{L4} R. P. Langlands, \textit{Reflexions on receiving the Shaw prize}, \url{http://publications.ias.edu/rpl/}.
\bibitem{L5} R. P. Langlands, \textit{Singularités et transferts}, en préparation.
\bibitem{Lm1} G. Laumon, \textit{Correspondance de Langlands géométrique pour les corps de fonctions}, Duke Math. J. 54 (1987), no. 2, p. 309--359.
\bibitem{Lm2} G. Laumon, \textit{Transformation de Fourier, constantes d'équations fonctionnelles et conjecture de Weil}, Inst. Hautes Études Sci. Publ. Math. No. 65 (1987), p. 131--210.
\bibitem{LMB} G. Laumon et L. Moret-Bailly, \textit{Champs algébriques}, Springer-Verlag, Berlin (2000).
\bibitem{LO} Yves Laszlo et Martin Olsson, \textit{The six operations for sheaves on Artin stacks. I. Finite coefficients}, Publ. Math. Inst. Hautes Études Sci. No. 107 (2008), 109--168; \textit{II. Adic coefficients}, Publ. Math. Inst. Hautes Études Sci. No. 107 (2008), 169--210.
\bibitem{MV} I. Mirković et K. Vilonen, \textit{Geometric Langlands duality and representations of algebraic groups over commutative rings}, Ann. of Math. (2) 166 (2007), no. 1, 95--143; \textit{Perverse sheaves on affine Grassmannians and Langlands duality}, Math. Res. Lett. 7 (2000), no. 1, 13--24.
\bibitem{N} B. C. Ngô, \textit{Le lemme fondamental pour les algèbres de Lie}, Publ. Math. Inst. Hautes Études Sci. No. 111 (2010), 1--169.
\bibitem{O1} T. Ono, \textit{Arithmetic of algebraic tori}, Ann. of Math., Vol. 74 (1961), p. 101--139.
\bibitem{O2} T. Ono, \textit{On the Tamagawa number of algebraic tori}, Ann. of Math., Vol. 78 (1963), p. 47--73.
\bibitem{O3} T. Ono, \textit{On the relative theory of Tamagawa numbers}, Ann. of Math., Vol. 82 (1965), p. 88--111.
\bibitem{R} M. Raynaud, \textit{Caractéristique d'Euler-Poincaré d'un faisceau et cohomologie des variétés abéliennes}, Séminaire Bourbaki, Vol. 9, Exp. No. 286, 129--147.
\bibitem{Sh} Joseph Shalika, \textit{A theorem on semi-simple p-adic groups}, Ann. of Math., Vol. 95 (1972), p. 226--242.
\bibitem{St1} Robert Steinberg, \textit{Regular elements of semisimple algebraic groups}, Pub. Sci. de l'IHES 25 (1965), p. 49--80.
\bibitem{St2} Robert Steinberg, \textit{Abstract homomorphisms of simple algebraic groups}, Sém. Bourbaki vol. 1972/73, p. 307--326.
\bibitem{V} V. S. Varadarajan, \textit{Harmonic analysis on real reductive groups}, SLN 576, Springer Verlag (1977).
\bibitem{WS} Wang Song, \url{http://www.mcm.ac.cn/people/ws.aspx}.
\bibitem{W1} A. Weil, \textit{Adeles and algebraic groups}, Institute for Advanced Study, Birkhäuser (1982).
\bibitem{W2} A. Weil, \textit{Basic number theory}, Springer-Verlag (1991).
\end{thebibliography}
